% from ega0-22-fr.tex
\oldpage[0\textsubscript{IV-1}]{193}

du complété $\widehat B$, tel que $\mathfrak q\cap B=0$, l'anneau local
$\widehat B_\mathfrak q$ est une $K$-algèbre formellement lisse pour sa
topologie $\mathfrak q$-adique, donc un anneau géométriquement régulier sur
$K$.
\end{corollary}

\begin{proof}
En effet, il résulte de (19.8.9) qu'il existe un sous-anneau $A_0$ de $A$
qui est un anneau de séries formelles sur un corps de caractéristique $p>0$
ou sur un anneau de Cohen (on rappelle que les anneaux de Cohen ont un corps
des fractions de caractéristique $0$), et est tel que $A$ soit une
$A_0$-algèbre finie. Or, on sait que l'anneau $A_0$ est régulier (17.3.8) et
vérifie l'une des hypothèses (i), (ii) de (22.3.2), en vertu de (21.8.8)~;
il suffit donc d'appliquer (22.3.2), en remplaçant $B$ par $A$ et $A$ par
$A_0$.
\end{proof}

\subsection{Résultats préliminaires sur les extensions finies d'anneaux locaux dont l'idéal maximal est de carré nul}
\label{subsection:0.22.4-fr}
\label{0.22.4-fr}

\begin{proposition}[22.4.1]
\label{0.22.4.1-fr}
Soient $A$ un anneau local dont l'idéal maximal $\mathfrak m$ est de carré
nul, $K=A/\mathfrak m$ son corps résiduel~; désignons par $V$ l'idéal
$\mathfrak m$ considéré comme $K$-espace vectoriel. Soient $T_i$
$(1\leqslant i\leqslant r)$ des indéterminées, et pour chaque $i$, soit
$F_i$ un polynôme unitaire de $A[T_i]$ de degré $d_i$~; désignons par $A'$ le
quotient de l'anneau de polynômes $A[T_1,\ldots,T_r]$ par l'idéal
$\mathfrak b$ engendré par les $r$ polynômes $F_i$. Supposons que $A'$ soit
un anneau local, et soient $\mathfrak m'$ son idéal maximal,
$K'=A'/\mathfrak m'$ son corps résiduel. Alors~:
\begin{enumerate}
  \item[(i)] Si l'on pose $d=\prod_{i=1}^r d_i$, $A'$ est un $A$-module
  libre de rang $d$, et l'on a $[K':K]\leqslant d$.
  \item[(ii)] Pour que $[K':K]=d$, il faut et il suffit que
  $A'\otimes_AK=A'/\mathfrak mA'$ soit un corps (nécessairement isomorphe à
  $K'$), autrement dit, que $\mathfrak m'=\mathfrak mA'$. On a dans ce cas
  $\mathfrak m'^2=0$, et si $V'$ désigne l'idéal $\mathfrak m'$ considéré
  comme $K'$-espace vectoriel, l'homomorphisme canonique
  $V\otimes_KK'\to V'$ est bijectif (et par suite
  $\operatorname{rg}_{K'}V'=\operatorname{rg}_K V$).
\end{enumerate}
\end{proposition}

\begin{proof}
Sans supposer $A'$ local, il est évident, par division euclidienne et
récurrence sur $r$, que si $t_i$ est la classe de $T_i$ mod.~$\mathfrak b$,
les monômes
$M_\nu=t_1^{\nu_1}t_2^{\nu_2}\cdots t_r^{\nu_r}$, où
$0\leqslant\nu_i\leqslant d_i-1$, forment une base du $A$-module $A'$. Si
$A'$ est supposé local, $K'$ est un quotient de
$A'/\mathfrak mA'=A'\otimes_AK$, donc
$[K':K]\leqslant[A'\otimes_AK:K]=d$, et il ne peut y avoir égalité que si
$\mathfrak m'=\mathfrak mA'$, ce qui démontre (i) et la première assertion
de (ii). D'autre part, il est clair que lorsque
$\mathfrak m'=\mathfrak mA'$, on a $\mathfrak m'^2=0$, et
$V\otimes_KK'=\mathfrak m\otimes_AA'=\mathfrak mA'=V'$ puisque $A'$ est
un $A$-module libre.
\end{proof}

\begin{remark}[22.4.2]
\label{0.22.4.2-fr}
Lorsque $A$ est artinien (autrement dit $\operatorname{rg}_K(V)$ fini),
l'hypothèse que $A'$ est local entraîne toujours que
$\operatorname{rg}_{K'}(\mathfrak m'/\mathfrak m'^2)
\geqslant\operatorname{rg}_K(V)$, comme nous le verrons plus loin (22.5.2).
On notera que ces deux rangs peuvent être égaux sans que l'on ait
$\mathfrak m'=\mathfrak mA'$.
\end{remark}

\begin{proposition}[22.4.3]
\label{0.22.4.3-fr}
Les hypothèses sur $A$ et les notations étant celles de (22.4.1), supposons
que les $F_i$ soient de la forme
\begin{equation}
\label{0.22.4.3.1-fr}
  F_i(T_i)=T_i^{d_i}+\sum_{1\leqslant k\leqslant d_i}
  c_{ik}T_i^{d_i-k},
  \tag{22.4.3.1}
\end{equation}
où $d_i\geqslant2$ pour tout $i$, et $c_{ik}\in\mathfrak m$
(«~polynômes d'Eisenstein~»). Alors~:
\begin{enumerate}
  \item[(i)] $A'$ est un anneau local, et son corps résiduel
  $K'=A'/\mathfrak m'$ est isomorphe à $K$.

  